\section{Uvod}

Digitalne platforme postale su svakodnevni prostor komunikacije, informisanja, zabave i društvene vidljivosti. Međutim, njihova uloga odavno nije jednostavno tehničko posredovanje sadržaja. Moderne platforme sadržaj rangiraju, preporučuju, testiraju i prilagođavaju s ciljem povećanja verovatnoće korisničke reakcije. U tom procesu pažnja korisnika postaje resurs koji se meri, predviđa i optimizuje.

Kroz rad je obrazložen stav da algoritamska manipulacija pažnjom ne nastaje kao izolovana posledica pojedinačnih (loših) dizajnerskih odluka, već kao rezultat spoja ekonomskog interesa, tehničke optimizacije, merenja korisničkog ponašanja i psiholoških mehanizama zadržavanja. Drugim rečima, problem nije samo u tome što korisnici provode mnogo vremena na platformama, već u tome što su platforme sistemski i ciljano podešene da to vreme produžavaju, analiziraju i pretvaraju u poslovnu vrednost.

Rad je organizovan kroz nekoliko povezanih celina. Najpre se razmatra širi društveni okvir ekonomije pažnje i nadzornog kapitalizma, kako bi se pokazalo zašto pažnja i podaci imaju centralnu ulogu u poslovnom modelu digitalnih platformi. Zatim se analiziraju centralne tehničke metode maksimizacije engagementa (tzv. collaborative filtering, reinforcement learning i A/B testiranje), kao i metrike kojima se pažnja kvantifikuje (poput klikova, vremena zadržavanja i retention-a). Na kraju se razmatraju psihološki mehanizmi i posledice po korisnike, uključujući promenljivo nagrađivanje, uklanjanje tačke prekida, socijalnu validaciju, FOMO i digitalnu adikciju.

Kroz rad se ne tvrdi da su sve digitalne platforme same po sebi štetne, niti se osporava korisnost personalizacije, preporuka ili analitike. Cilj rada jeste da ukaže na to da iste tehnike koje mogu poboljšati korisničko iskustvo, u poslovnom modelu zasnovanom na pažnji, mogu postati sredstva za slabljenje korisničke autonomije, manipulaciju ponašanjem i indukovanje digitalne adikcije.